\documentclass[letterpaper,11pt]{article}
\usepackage{resume_style}
\begin{document}

%---------- HEADER ----------
\begin{center}
  {\Huge \scshape Vedaang Chopra} \\ \vspace{2pt}
  \small
  \href{mailto:vedaangchopra@gatech.edu}{\underline{vedaangchopra@gatech.edu}} \;|\;
  \underline{+1 (404) 740-9905} \;|\;
  Atlanta, GA \;|\;
  \href{https://linkedin.com/in/vedaang-chopra}{\underline{LinkedIn}} \;|\;
  \href{https://github.com/Vedaang-Chopra}{\underline{GitHub}} \;|\;
  \href{https://vedaangchopra.live/}{\underline{Website}}
  \\ \vspace{3pt}
  \small
  MS CS (ML) at Georgia Tech \;|\; Advised by \textbf{Dr.\ Matthew Gombolay}, CORE Robotics Lab \\
  \textit{Research:} Coding agents with execution-grounded verification -- Tool-use planning \& multi-agent orchestration -- Evaluation harnesses that expose model failures
\end{center}
\vspace{-6pt}

%---------- RESEARCH EXPERIENCE ----------
\section{Research Experience}
\resumeSubHeadingListStart
\resumeSubheading
{Georgia Institute of Technology -- CORE Robotics Lab @ Siemens}{Atlanta, GA}
{Graduate Research Assistant \;|\; Advisor: \textbf{Dr.\ Matthew Gombolay}}{Jan 2026 -- Present}
\resumeItemListStart
  \resumeItem{Building a closed-loop coding-agent pipeline for parametric CAD code generation:
  a frozen LLM writes CadQuery programs, execution feedback (compile validity, geometric
  correctness via Chamfer/Hausdorff on STL/B-Rep) drives a LangGraph-orchestrated debug-and-repair
  loop -- framing synthesis as inference-time search with verifiable geometric reward; 16-module
  system with a 202-test automated evaluation harness.}

  \resumeItem{Designing graders and failure diagnostics for generated code: AST/code-graph structural
  analysis, failure classification, and repair-trajectory quality scoring -- turning messy qualitative
  behavior into concrete signals that steer iterative repair; VLM-based visual judgment complements
  programmatic metrics.}
\resumeItemListEnd

\resumeSubheading
{Georgia Institute of Technology -- CS 8903 Special Problems}{Atlanta, GA}
{Graduate Researcher}{Jan 2025 -- Present}
\resumeItemListStart
  \resumeItem{\textbf{SHASTRA} (Jan 2026--Present): Designing a trace-to-graph pipeline converting agent
  execution traces (GAIA, 94 sessions, 4,037 events, 442 graph nodes) into reusable task-composition
  graphs with dependency/control/conditional edges -- compositional workflow reuse for new long-horizon
  tasks under cost/latency constraints.}

  \resumeItem{\textbf{ATHENA} (Jan--May 2025): Built a 5-agent supervisor-worker multi-agent system via
  LangGraph Command routing with reflection-based critique and dynamic plan modification; structured
  outputs validated with Pydantic; +0.18 BLEU-4 absolute over single-agent baseline on 100 reference
  videos; ablation showed removing the critique agent cut fact coverage 27\%.}

  \resumeItem{\textbf{ARTEMIS} (Aug 2025--Present): Cost-aware VLM routing as dynamic tool selection --
  trained neural router picks the optimal model per request across 6+ backends served via vLLM with
  accuracy/cost/latency modes and SLA-aware load balancing.}
\resumeItemListEnd
\resumeSubHeadingListEnd

%---------- INDUSTRY EXPERIENCE ----------
\section{Industry Experience}
\resumeSubHeadingListStart
\resumeSubheading
{Fortinet Technologies Inc.\ (AIOps R\&D)}{Bengaluru, India}
{Software Development Engineer I \& II (ML/AI)}{Feb 2021 -- Jul 2025}
\resumeItemListStart
  \resumeItem{Led development of an agentic RAG diagnostics system from hackathon prototype to
  production deployment: LLM plans multi-step tool use over network telemetry, invokes APIs, and
  verifies hypotheses via structured function calling before concluding -- autonomous root-cause
  analysis reducing mean resolution time ~70\%.}

  \resumeItem{Scaled OpenSearch ingestion pipelines 50$\to$2,000 events/sec (40x) with Golang;
  deployed edge ML inference via ONNX Runtime (~40\% latency reduction); built SLA forecasting
  for 60+ classifiers across 4 categories.}

  \resumeItem{Patent (Filed): PCT/IN2022/058026 -- Unsupervised distributional thresholding for
  wireless connectivity anomaly detection; reduced manual troubleshooting ~75\%.}
\resumeItemListEnd
\resumeSubHeadingListEnd

%---------- TECHNICAL SKILLS ----------
\section{Technical Skills}
\footnotesize
\textbf{Agentic AI:} Coding Agents, LangChain, LangGraph,
  Agentic RAG, Tool/Function Calling, Multi-Agent Collaboration, Long-Horizon Task Decomposition,
  Structured Generation (Pydantic), Workflow Composition \\[1pt]
\textbf{Evals \& Verification:} Automated Evaluation Harnesses, Grader Design (Execution-Grounded),
  LLM-as-a-Judge, Ablation Studies, Model Evaluation \& Benchmarking, Failure Classification,
  Geometric/Programmatic Verifiers \\[1pt]
\textbf{AI / ML:} PyTorch, Hugging Face Transformers, Deep Learning, Transformers, NLP, Computer Vision,
  Scikit-learn, NumPy, Pandas \\[1pt]
\textbf{LLMs / VLMs / Serving:} vLLM, VLM Routing \& Evaluation, Cross-Modal Retrieval, CLIP, SBERT, Inference Optimization \\[1pt]
\textbf{Systems:} Python, Go, C/C++, SQL/PostgreSQL, Redis, OpenSearch/Elasticsearch, Docker, Linux, Git,
  FAISS, ONNX Runtime, FastAPI, Azure, W\&B, TensorBoard \\[1pt]
\textbf{Reinforcement Learning:} PPO, DQN, Reward Shaping, Curriculum Learning, Self-Play (academic projects)

%---------- EDUCATION ----------
\section{Education}
\resumeSubHeadingListStart
\resumeSubheading
{Georgia Institute of Technology}{Atlanta, GA}
{\textbf{M.S.\ in Computer Science} (Machine Learning Specialization), \textbf{GPA: 4.0/4.0}}{Aug 2024 -- Dec 2026}
\vspace{2pt}
\begin{itemize}[leftmargin=0.20in, topsep=2pt, itemsep=0pt]
  \item\footnotesize{\textbf{Coursework:} Deep Learning, Deep RL, ML Security, Large \& Vision Language Models, Agentic AI, Systems for AI}
\end{itemize}
\vspace{-2pt}
\resumeEducationCompact
{Maharaja Surajmal Institute of Technology}{New Delhi, India}
{B.Tech.\ in Information Technology, CGPA: 8.8/10.0}{Aug 2016 -- Aug 2020}
\resumeSubHeadingListEnd

\end{document}
